% !PRINTED 626
\weberpage{626}

The value $u=0$ or $v=0$ occurs only when
$\alpha=0$. In that case we have in fact in (13) the well-known
metacyclic equation $x^5+\beta=0$; this case, then, we
now set aside. The equation (20) can only have
two opposite roots when $\Delta=0$, that is,
when the equation (13) has equal roots. This too is
to be excluded, since we have assumed (13) to be irreducible.

There remains, then, only the question whether (20) can have
equal roots. In that case, together with (20), the derived
equation
\[
6u^5-20\alpha u^3+30\alpha^2u-\sqrt{\Delta}=0
\]
would have to be satisfied, and if we eliminate $\sqrt{\Delta}$ from this
and from (20),
\[
5u^6-15\alpha u^4+15\alpha^2u^2-5\alpha^3=5(u^2-\alpha)^3=0.
\]
We should therefore have to have $v=\alpha$, and hence by (22) $\beta=0$.
But then the equation of degree 5 would again be reducible, since
it has the factor $x$.

If one wishes, then, to test a given equation of the form (13)
with regard to its solvability, one has first to establish the
irreducibility, and then to investigate whether (21)
or (22) has a rational root.

If $\alpha$ and $\beta$ are integers, then a rational $v$ must also
be an integer, which is to be sought among the factors
of $25\alpha^6$.

If, e.\,g., $\alpha=5$, $\beta=5t$, then, if $t$ is an integer not
divisible by 5, $x^5+\alpha x+\beta$ is irreducible by \S.\,178, 3.
For $v$ one has to substitute powers of 5, with positive or
negative sign, but for no such number
can (22) be satisfied, because on the left-hand side the power of
5 does not become high enough.

\gesperrt{Thus no equation of the form}
$x^5+5x+5t=0$ \gesperrt{is metacyclic}.

If one wishes to determine metacyclic equations, one sets
in (22)
\[
\beta=\alpha\mu,\quad v=\alpha\lambda,
\]
whereby one obtains:
\weq{23}{\begin{aligned}
\alpha&=\dfrac{5^5\mu^4\lambda}{(\lambda-1)^4(\lambda^2+6\lambda+25)}\\
\beta&=\dfrac{5^5\mu^5\lambda}{(\lambda-1)^4(\lambda^2+6\lambda+25)},
\end{aligned}}
%CONT>>
